% =====================================================================
% Full derivation of IF(q_r) for ridge-penalized logistic regression
% in the MCD setting, with every object defined, the zero-mean check,
% and the practical algorithm.
%
% Self-contained preamble: compiles alone. All macros of the companion
% paper are redefined here, so the notation matches and nothing breaks
% if the section is pasted into the main file (delete the preamble then).
% =====================================================================

\documentclass[11pt]{article}
\usepackage[T1]{fontenc}
\usepackage[margin=1in]{geometry}
\usepackage{amsmath,amssymb,amsthm}
\usepackage{bm}
\usepackage{mathtools}
\usepackage{hyperref}
\usepackage[round]{natbib}
\usepackage{parskip}
\usepackage{xcolor}
\usepackage[most]{tcolorbox}
\usepackage{booktabs}
\usepackage{algorithm}
\usepackage{algpseudocode}

\newtheorem{proposition}{Proposition}
\newtheorem{lemma}{Lemma}
\newtheorem{corollary}{Corollary}
\newtheorem*{remark}{Remark}
\newtheorem*{assumption}{Assumption}

% --- macros of the companion paper -----------------------------------
\newcommand{\dd}{\,\mathrm{d}}
\newcommand{\Eb}{\mathbb{E}}
\newcommand{\Tset}{\mathcal{T}}
\newcommand{\Xset}{\mathcal{X}}
\newcommand{\IF}{\mathrm{IF}}
\newcommand{\Hb}{\mathcal{H}}
\newcommand{\Dz}{\mathcal{D}_Z}
% ---------------------------------------------------------------------

\newtcolorbox{keyresult}[1][]{
  colback=blue!4, colframe=blue!45!black, boxrule=0.4pt,
  arc=1mm, left=6pt, right=6pt, top=4pt, bottom=4pt, #1
}

\begin{document}

\section{The influence function of a ridge-penalized logistic classifier}
\label{sec:if-logistic}

This section derives $\IF(o;q_r,F)$ in full for a ridge-penalized logistic
classifier, defines every object it uses, verifies the result against the zero-mean
property that every influence function must satisfy, and gives the algorithm that
computes it. The same formulas serve the multilayer perceptron and CatBoost once
their representation is frozen, since both are logistic models on a learned feature
map; only $\phi$ changes.

\subsection{Objects and definitions}
\label{sec:objects}

\paragraph{The augmented sample.}
The MCD construction turns conditional density estimation into binary classification.
The augmented observation is $W=(T,X,Y,Z)$: with probability $r\in(0,1)$ the triple
$(T,X,Y)$ is a real observation, labeled $Z=1$; with probability $1-r$ it is
contrastive, $(T,X,Y')$ with $Y'$ drawn from the marginal $f_Y$ independently of
$(T,X)$, labeled $Z=0$. Write
\[
P_1(F) := F_{TXY} \quad\text{(the real law)},
\qquad
P_0(F) := F_{TX}\otimes F_Y \quad\text{(the contrastive law)},
\]
and let $\Eb_{P_r}:=r\,\Eb_{P_1}+(1-r)\,\Eb_{P_0}$ denote expectation under the
augmented mixture. The structural fact that drives the whole derivation is that
$P_0(F)$ is itself a functional of $F$, built from two of its marginals: perturbing
$F$ moves the contrastive class as well as the real one.

\paragraph{The link.}
The logistic sigmoid $\sigma:\mathbb{R}\to(0,1)$ is
\[
\sigma(u)=\frac{1}{1+e^{-u}},
\qquad
\sigma'(u)=\sigma(u)\big(1-\sigma(u)\big).
\]

\paragraph{The feature map.}
$\phi:\Tset\times\Xset\times\mathbb{R}\to\mathbb{R}^k$ maps a raw input to a vector of
$k$ numbers. For an ordinary logistic model it is a fixed basis, for instance
$\phi(t,x,y)=(1,t^\top,x^\top,y,\text{interactions})^\top$; for a frozen network it is
the penultimate-layer activation; for a frozen boosted ensemble it is the
concatenated leaf-indicator vector. What matters here is only that $\phi$ is fixed and
known, not estimated from the fold on which we differentiate.

\paragraph{The model, the loss, the criterion.}
With the logit $s:=\beta^\top\phi(w)$,
\[
q_\beta(w) := \sigma\big(\beta^\top\phi(w)\big),
\qquad
\ell(z,q) := -z\log q-(1-z)\log(1-q),
\]
the cross-entropy. The population criterion and its minimizer, the estimand, are
\begin{equation}
M_\lambda(\beta;F)
= r\,\Eb_{P_1}\big[\ell(1,q_\beta)\big]
+ (1-r)\,\Eb_{P_0}\big[\ell(0,q_\beta)\big]
+ \lambda\|\beta\|_2^2,
\qquad
\beta_\lambda(F)=\arg\min_{\beta\in\mathbb{R}^k} M_\lambda(\beta;F).
\label{eq:criterion}
\end{equation}
The classical unpenalized fit is the case $\lambda=0$.

\subsection{Gradient and Hessian of the loss}
\label{sec:grads}

Differentiating $\ell$ in its second argument,
\[
\frac{\partial\ell}{\partial q}
= -\frac{z}{q}+\frac{1-z}{1-q}
= \frac{-z(1-q)+(1-z)q}{q(1-q)}
= \frac{q-z}{q(1-q)},
\]
and differentiating the model in $\beta$, using $\sigma'=\sigma(1-\sigma)$,
\begin{equation}
\nabla_\beta q_\beta = \sigma'(s)\,\nabla_\beta s = q_\beta(1-q_\beta)\,\phi.
\label{eq:grad-q}
\end{equation}
The chain rule gives the classical cancellation of the canonical link,
\begin{equation}
\nabla_\beta\,\ell\big(z,q_\beta\big)
= \frac{q_\beta-z}{q_\beta(1-q_\beta)}\cdot q_\beta(1-q_\beta)\,\phi
= \big(q_\beta-z\big)\,\phi,
\label{eq:grad-loss}
\end{equation}
so the two per-label gradients are
\begin{equation}
u_\beta := \nabla_\beta\ell(1,q_\beta) = (q_\beta-1)\,\phi,
\qquad
v_\beta := \nabla_\beta\ell(0,q_\beta) = q_\beta\,\phi.
\label{eq:uv}
\end{equation}
Differentiating once more, and using~\eqref{eq:grad-q},
\[
\nabla_\beta u_\beta = \phi\,(\nabla_\beta q_\beta)^\top = q_\beta(1-q_\beta)\,\phi\phi^\top
= \nabla_\beta v_\beta,
\]
so both labels contribute the same second derivative and the two expectations merge
into a single one under the mixture:
\begin{equation}
\Hb_\lambda := \nabla^2_\beta M_\lambda
= \Eb_{P_r}\big[q_\beta(1-q_\beta)\,\phi\phi^\top\big] + 2\lambda I_k.
\label{eq:hessian}
\end{equation}
Each $\phi\phi^\top$ is positive semidefinite and the weight $q(1-q)$ is strictly
positive, so the expectation is positive semidefinite; adding $2\lambda I_k$ makes
$\Hb_\lambda$ positive definite with all eigenvalues at least $2\lambda>0$. It is
therefore invertible and well conditioned, uniformly in the sample. For $\lambda=0$
invertibility requires only that $\phi$ exhibit no perfect collinearity on the
support.

The first-order condition characterizing $\beta_\lambda(F)$ is
\begin{equation}
G_\lambda(\beta;F) := \nabla_\beta M_\lambda
= r\,\Eb_{P_1}[u_\beta] + (1-r)\,\Eb_{P_0}[v_\beta] + 2\lambda\beta = 0,
\label{eq:foc}
\end{equation}
which we use twice below, once to simplify and once to check.

\subsection{Perturbing the augmented law}
\label{sec:perturb}

Let $F_\varepsilon=(1-\varepsilon)F+\varepsilon\delta_o$ with $o=(y_o,t_o,x_o)$. By
definition $\beta_\lambda(F_\varepsilon)$ solves
$G_\lambda(\beta_\lambda(F_\varepsilon);F_\varepsilon)=0$ identically in $\varepsilon$.
Differentiating this identity at $\varepsilon=0$ and applying the implicit function
theorem, legitimate because $\Hb_\lambda$ is nonsingular,
\begin{equation}
\Hb_\lambda\cdot\IF(o;\beta_\lambda,F) + D(o;\beta) = 0,
\qquad
D(o;\beta):=\frac{\partial}{\partial\varepsilon}G_\lambda(\beta;F_\varepsilon)\Big|_{\varepsilon=0},
\label{eq:ift}
\end{equation}
so that $\IF(o;\beta_\lambda,F)=-\Hb_\lambda^{-1}D(o;\beta)$. It remains to compute
$D$, that is, to differentiate the two expectations in~\eqref{eq:foc} at frozen
$\beta$.

\paragraph{The real channel.}
Since $P_1(F_\varepsilon)=F_\varepsilon$ and
$\Eb_{F_\varepsilon}[h]=(1-\varepsilon)\Eb_F[h]+\varepsilon\,h(o)$ is affine in
$\varepsilon$,
\begin{equation}
\frac{\partial}{\partial\varepsilon}\Eb_{P_1(F_\varepsilon)}[u_\beta]\Big|_{0}
= u_\beta(t_o,x_o,y_o) - \Eb_{P_1}[u_\beta].
\label{eq:channel1}
\end{equation}

\paragraph{The contrastive channels.}
The perturbed marginals are
$F_{TX,\varepsilon}=(1-\varepsilon)F_{TX}+\varepsilon\delta_{(t_o,x_o)}$ and
$F_{Y,\varepsilon}=(1-\varepsilon)F_Y+\varepsilon\delta_{y_o}$, and
$P_0(F_\varepsilon)=F_{TX,\varepsilon}\otimes F_{Y,\varepsilon}$. Expanding the double
integral in full,
\begin{equation}
\begin{aligned}
\Eb_{P_0(F_\varepsilon)}[v_\beta]
&= \iint v_\beta(t,x,y')\dd F_{TX,\varepsilon}(t,x)\dd F_{Y,\varepsilon}(y')\\[2pt]
&= (1-\varepsilon)^2\,\Eb_{P_0}[v_\beta]
\;+\;(1-\varepsilon)\varepsilon\,\Eb_{F_{TX}}\big[v_\beta(T,X,y_o)\big]\\
&\qquad+\;\varepsilon(1-\varepsilon)\,\Eb_{F_{Y}}\big[v_\beta(t_o,x_o,Y')\big]
\;+\;\varepsilon^2\,v_\beta(t_o,x_o,y_o),
\end{aligned}
\label{eq:p0-expanded}
\end{equation}
the four terms corresponding to neither, one, the other, or both factors being
contaminated. Differentiating at $\varepsilon=0$, the quadratic term vanishes, the
$(1-\varepsilon)^2$ term contributes $-2$, and each cross term contributes $+1$:
\begin{equation}
\frac{\partial}{\partial\varepsilon}\Eb_{P_0(F_\varepsilon)}[v_\beta]\Big|_{0}
= \Eb_{F_Y}\big[v_\beta(t_o,x_o,Y')\big]
+ \Eb_{F_{TX}}\big[v_\beta(T,X,y_o)\big]
- 2\,\Eb_{P_0}[v_\beta].
\label{eq:channel0}
\end{equation}
The factor $-2$ is what splits the contrastive derivative into two channels, one per
factor of the product measure. Writing them separately,
\begin{equation}
\begin{aligned}
D_1 &:= r\Big(u_\beta(t_o,x_o,y_o)-\Eb_{P_1}[u_\beta]\Big),\\
D_0^{TX} &:= (1-r)\Big(\Eb_{F_Y}\big[v_\beta(t_o,x_o,Y')\big]-\Eb_{P_0}[v_\beta]\Big),\\
D_0^{Y} &:= (1-r)\Big(\Eb_{F_{TX}}\big[v_\beta(T,X,y_o)\big]-\Eb_{P_0}[v_\beta]\Big),
\end{aligned}
\label{eq:channels}
\end{equation}
so that $D=D_1+D_0^{TX}+D_0^{Y}$; the penalty $2\lambda\beta$ does not depend on $F$
and contributes nothing, and the mixing weight $r$ is a design constant. The first
contrastive channel freezes the contaminated pair $(t_o,x_o)$ and averages over
$Y'\sim f_Y$; the second freezes the contaminated outcome $y_o$ and averages over
$(T,X)\sim f_{TX}$. These two channels are the price of manufacturing the negative
class from the data, and they are absent from the textbook sandwich formula, which
conditions on a fixed training law.

\subsection{The closed form}
\label{sec:closed-form}

Introduce the three vectors of $\mathbb{R}^k$
\begin{equation}
b_{TX}(t_o,x_o) := \Eb_{Y'\sim f_Y}\big[q_\beta\phi\,(t_o,x_o,Y')\big],
\qquad
b_{Y}(y_o) := \Eb_{(T,X)\sim f_{TX}}\big[q_\beta\phi\,(T,X,y_o)\big],
\qquad
b_0 := \Eb_{P_0}\big[q_\beta\phi\big],
\label{eq:b-vectors}
\end{equation}
so that, by~\eqref{eq:uv}, $\Eb_{F_Y}[v_\beta(t_o,x_o,Y')]=b_{TX}(t_o,x_o)$,
$\Eb_{F_{TX}}[v_\beta(T,X,y_o)]=b_Y(y_o)$ and $\Eb_{P_0}[v_\beta]=b_0$. Summing the
three channels of~\eqref{eq:channels},
\[
D(o;\beta) = r\,u_\beta(o) + (1-r)\big(b_{TX}(t_o,x_o)+b_Y(y_o)\big)
- \Big(r\,\Eb_{P_1}[u_\beta] + 2(1-r)\,b_0\Big),
\]
the last bracket collecting the three centering terms, one from $D_1$ and two
identical ones from the contrastive channels. The first-order condition~\eqref{eq:foc}
gives $r\,\Eb_{P_1}[u_\beta] = -2\lambda\beta_\lambda-(1-r)b_0$, whence
\[
-\Big(r\,\Eb_{P_1}[u_\beta]+2(1-r)b_0\Big)
= 2\lambda\beta_\lambda+(1-r)b_0-2(1-r)b_0
= 2\lambda\beta_\lambda-(1-r)\,b_0 .
\]
Note that a residual $-(1-r)b_0$ survives: the first-order condition removes one of
the two contrastive centering terms but not both, and dropping it would be an error.
This gives the result.

\begin{keyresult}
For ridge-penalized logistic regression with $\lambda\ge0$, at
$\beta=\beta_\lambda(F)$ and $o=(y_o,t_o,x_o)$,
\begin{equation}
\IF\big(o;\beta_\lambda,F\big)
= -\,\Hb_\lambda^{-1}\,D(o;\beta),
\qquad
D(o;\beta) = r\big(q_\beta(o)-1\big)\phi(o)
+ (1-r)\Big(b_{TX}(t_o,x_o)+b_Y(y_o)-b_0\Big)
+ 2\lambda\beta_\lambda,
\label{eq:if-beta}
\end{equation}
with $\Hb_\lambda$ as in~\eqref{eq:hessian} and $b_{TX},b_Y,b_0$ as
in~\eqref{eq:b-vectors}. By the chain rule and~\eqref{eq:grad-q}, for any fixed
evaluation point,
\begin{equation}
\IF\big(o;q_r(t,x,y),F\big)
= \nabla_\beta q_\beta(t,x,y)^\top\IF\big(o;\beta_\lambda,F\big)
= q_\beta(1-q_\beta)\,\phi(t,x,y)^\top\,\IF\big(o;\beta_\lambda,F\big).
\label{eq:if-qr}
\end{equation}
\end{keyresult}

Setting $\lambda=0$ recovers the classical logistic case. Dropping in addition the two
contrastive channels, which is what fixing the contrast law by design would do as in
noise-contrastive estimation \citep{gutmann2010nce}, leaves
$-\Hb_0^{-1}(q_\beta(o)-z_o)\phi(o)$, the textbook sandwich formula.

\subsection{Verification: the zero-mean property}
\label{sec:check}

Every influence function integrates to zero under the law at which it is taken. This
is the sharpest available check on~\eqref{eq:if-beta}, and it fails if the residual
$-b_0$ is omitted.

\begin{lemma}[Zero mean]
\label{lem:zero-mean}
$\Eb_{o\sim F}\big[\IF(o;\beta_\lambda,F)\big]=0$, equivalently
$\Eb_{o\sim F}[D(o;\beta)]=0$.
\end{lemma}

\begin{proof}
Take the expectation of $D$ in~\eqref{eq:if-beta} over $o\sim F$, term by term. For the
real channel, $o$ is distributed as $P_1$, so
$\Eb_o[r(q_\beta(o)-1)\phi(o)]=r\,\Eb_{P_1}[u_\beta]$. For the first contrastive
channel, the contaminated pair $(t_o,x_o)$ is distributed as $F_{TX}$, so by Fubini and
the independence built into $P_0$,
\[
\Eb_o\big[b_{TX}(t_o,x_o)\big]
= \Eb_{(T,X)\sim F_{TX}}\Eb_{Y'\sim F_Y}\big[v_\beta(T,X,Y')\big]
= \Eb_{P_0}[v_\beta] = b_0,
\]
and symmetrically $\Eb_o[b_Y(y_o)]=b_0$, since $y_o$ is distributed as $F_Y$. The
vectors $b_0$ and $2\lambda\beta_\lambda$ are constants. Summing,
\[
\Eb_o[D] = r\,\Eb_{P_1}[u_\beta] + (1-r)\big(b_0+b_0-b_0\big) + 2\lambda\beta_\lambda
= \underbrace{r\,\Eb_{P_1}[u_\beta]+(1-r)\,\Eb_{P_0}[v_\beta]}_{=-2\lambda\beta_\lambda\ \text{by}~\eqref{eq:foc}}
+\,2\lambda\beta_\lambda = 0. \qedhere
\]
\end{proof}

Without the residual $-b_0$ the sum would leave $(1-r)b_0\neq0$, so the term is not
optional. Lemma~\ref{lem:zero-mean} also yields a free numerical diagnostic: the
empirical average of $\widehat D(O_i)$ over the evaluation fold must be close to zero,
and a departure signals an error in $\hat b_0$, in one of the channels, or in the
empirical first-order condition.

\subsection{Plugging into the interventional density}
\label{sec:plug}

The MCD-native correction requires the classifier's influence function only through
\[
\Eb_X\!\left[\frac{\IF\big(o;q_r(t,X,y),F\big)}{\big(1-q_r(t,X,y)\big)^2}\right].
\]
Since $\IF(o;\beta_\lambda,F)$ does not depend on $X$, it factors out of the
expectation, and by~\eqref{eq:if-qr} together with
$q(1-q)/(1-q)^2=q/(1-q)=g$,
\begin{equation}
\Eb_X\!\left[\frac{\IF\big(o;q_r(t,X,y),F\big)}{\big(1-q_r(t,X,y)\big)^2}\right]
= a_\lambda(t,y)^\top\,\IF\big(o;\beta_\lambda,F\big),
\qquad
a_\lambda(t,y) := \Eb_X\big[g(t,X,y)\,\phi(t,X,y)\big],
\label{eq:a-lambda}
\end{equation}
the odds-weighted average feature vector. Substituting into the MCD-native influence
function of the target gives the one-step estimator, in which no generalized
propensity score and no kernel in $T$ appear:
\begin{equation}
\hat p_t(y)
= \frac1n\sum_{i=1}^n\Big\{
\hat f_Y(y)\,\hat c\Big[\hat g(t,X_i,y)-\bar g(t,y)-\hat\xi(t,y)^\top\widehat D(O_i)\Big]
+ \hat m(t,y)\big(K_{h_Y}(Y_i-y)-\hat f_Y(y)\big)\Big\},
\label{eq:final-estimator}
\end{equation}
with $\bar g(t,y):=\frac1n\sum_j\hat g(t,X_j,y)$, $\hat m(t,y):=\hat c\,\bar g(t,y)$,
$\hat c=(1-r)/r$, and
\begin{equation}
\xi(t,y) := \Hb_\lambda^{-1}a_\lambda(t,y),
\qquad\text{so that}\qquad
a_\lambda(t,y)^\top\IF(o;\beta_\lambda,F) = -\,\xi(t,y)^\top D(o;\beta).
\label{eq:xi}
\end{equation}
The reparametrization~\eqref{eq:xi} matters computationally: it moves the linear solve
from the observations to the evaluation points, of which there are far fewer.

\subsection{Algorithm}
\label{sec:algo}

\paragraph{Building the augmented sample.}
For each $i$, keep the real quadruple $(T_i,X_i,Y_i,Z{=}1)$ and create $\kappa$
contrastive quadruples $(T_i,X_i,Y'_{i\ell},Z{=}0)$ with $Y'_{i\ell}$ drawn uniformly
from the observed outcomes $\{Y_1,\dots,Y_n\}$, which samples $f_Y$ by permutation.
Then $r=1/(1+\kappa)$ and $c=(1-r)/r=\kappa$; the choice $\kappa=1$ gives $r=1/2$ and
$c=1$. Note that $\Dz$ is a sample from $P_r$ and its contrastive half is a sample
from $P_0$, so both $\Hb_\lambda$ and $b_0$ are ordinary averages over $\Dz$, at cost
$O(N)$ rather than $O(n^2)$.

\begin{algorithm}[htbp]
\caption{One-step MCD estimator with a ridge logistic classifier}
\label{alg:main}
\begin{algorithmic}[1]
\State Split the sample into folds $\mathcal{A}$ (fitting) and $\mathcal{B}$ (evaluation); swap and average at the end.
\State On $\mathcal{A}$: build $\Dz$ as above and fit $\hat\beta_\lambda$ by ridge-penalized cross-entropy.
\State On $\mathcal{A}$: form $\hat\Hb_\lambda=\frac1N\sum_{w\in\Dz}\hat q(w)(1-\hat q(w))\phi(w)\phi(w)^\top+2\lambda I_k$ and factorize it once by Cholesky. \Comment{$O(Nk^2+k^3)$}
\State On $\mathcal{A}$: form $\hat b_0=\frac{1}{N_0}\sum_{w\in\text{contrastive half}}\hat q(w)\phi(w)$. \Comment{$O(N_0k)$}
\For{each evaluation point $(t,y)$}
  \State $\hat a_\lambda(t,y)\gets\frac1n\sum_j\hat g(t,X_j,y)\phi(t,X_j,y)$ with $\hat g=\hat q/(1-\hat q)$.
  \State Solve $\hat\xi(t,y)\gets\hat\Hb_\lambda^{-1}\hat a_\lambda(t,y)$ by back-substitution. \Comment{$O(k^2)$}
\EndFor
\For{each $O_i$ in $\mathcal{B}$}
  \State $\hat b_{TX}(T_i,X_i)\gets\frac1S\sum_{j\in\mathcal{S}}\hat q(T_i,X_i,Y_j)\phi(T_i,X_i,Y_j)$.
  \State $\hat b_{Y}(Y_i)\gets\frac1S\sum_{j\in\mathcal{S}}\hat q(T_j,X_j,Y_i)\phi(T_j,X_j,Y_i)$.
  \State $\widehat D(O_i)\gets r(\hat q(O_i)-1)\phi(O_i)+(1-r)\big(\hat b_{TX}(T_i,X_i)+\hat b_Y(Y_i)-\hat b_0\big)+2\lambda\hat\beta_\lambda$.
\EndFor
\State Check $\big\|\frac1n\sum_i\widehat D(O_i)\big\|\approx0$ (Lemma~\ref{lem:zero-mean}).
\State Return $\hat p_t(y)$ from~\eqref{eq:final-estimator}, each $i$ costing one inner product $\hat\xi(t,y)^\top\widehat D(O_i)$.
\end{algorithmic}
\end{algorithm}

\paragraph{Cost.}
The two contrastive channels are the only quadratic step in principle, since each
requires an average over the marginals for every $o$. Subsampling the inner average
over $S$ draws, with $S$ of the order of a few hundred, reduces this to $O(nSk)$; the
resulting Monte Carlo error is $O(S^{-1/2})$ and is absorbed, since these vectors enter
a correction that is itself averaged over $i$. Everything else is linear in $n$ apart
from the single Cholesky factorization, $O(k^3)$, and one back-substitution per
evaluation point, $O(k^2)$.

\paragraph{On the choice of $\lambda$.}
At fixed $\lambda>0$ the estimand is the ridge population minimizer $\beta_\lambda(F)$
and~\eqref{eq:if-beta} is exact for it; the term $2\lambda I_k$ inside $\Hb_\lambda$
bounds the influence, the local robustness of penalized M-estimators
\citep{avella2017}. At $\lambda_n=o(n^{-1/2})$ the penalty is first-order negligible
and~\eqref{eq:if-beta} reduces to the classical influence function. In both regimes the
closed form is used as written, and the only bandwidth in the whole procedure is $h_Y$,
carried by the marginal correction in~\eqref{eq:final-estimator}.

\bibliographystyle{plainnat}
\bibliography{references}

\end{document}